\begin{proposition}[an explicit nonzero local cusp restriction]
\hypertarget{tag-03IN}{}\label{wbII:prop:local-cusp-picard-obstruction}
\StatementRecord{wbII:prop:local-cusp-picard-obstruction}{proved}{03IN}{\SeriesRef{wbII:eq:W:5.17t}, \SeriesRef{wbII:eq:W:5.17ad}, \SeriesRef{wbII:eq:W:5.17u}, \SeriesRef{wbII:eq:W:5.17v}, \SeriesRef{wbII:eq:W:5.17ai}, \SeriesRef{wbII:eq:W:5.17ag}, \SeriesRef{wbII:eq:W:5.17aj}, \SeriesRef{wbII:eq:W:5.17ak}, \SeriesRef{wbII:eq:W:5.17al}, \SeriesRef{wbII:eq:W:5.17am}, \SeriesRef{wbII:eq:W:5.17an}, \SeriesRef{wbII:eq:W:5.17ao}, \SeriesRef{wbII:eq:W:5.17ae}, \SeriesRef{wbII:eq:W:5.17ax}}

\ Let
\[
 H=\begin{pmatrix}1&0\\0&-6\end{pmatrix},
 \qquad
 P=B_0^{\mathsf t}H=\begin{pmatrix}0&6\\1&0\end{pmatrix},
 \qquad
 Q(v)=v^{\mathsf t}Hv=v_1^2-6v_2^2.
\tag{W.65}\hypertarget{tag-03HF}{}\label{wbII:eq:W:5.17x}
\]
Then
\[
 \mathcal D_Q=\sum_{v\in\mathbb Z^2}Q(v)E_v
\tag{W.66}\hypertarget{tag-03HG}{}\label{wbII:eq:W:5.17y}
\]
is a locally finite Cartier divisor on \(Y_\epsilon\).
The line bundle \(\mathcal O_{Y_\epsilon}(\mathcal D_Q)\) has an explicit
\((\bar\Lambda,\Psi)\)-equivariant structure and descends to a holomorphic
line bundle \(\mathcal L_Q\) on \(N_0\).  If
\[
 (a_1,a_2,a_3,a_4,a_5,a_6)
 =(e_1,e_2,e_2-e_1,-e_1,-e_2,e_1-e_2),
 \qquad C_i=E_0\cap E_{a_i},
\]
then
\[
 \nu^*(\mathcal L_Q|_W)
 \simeq
 \mathcal O_{E_0}\!\left(\sum_{i=1}^6Q(a_i)C_i\right),
\qquad
 \left(\deg_{C_i}\nu^*(\mathcal L_Q|_W)\right)_{i=1}^6
 =(-12,2,0,-12,2,0).
\tag{W.67}\hypertarget{tag-03HH}{}\label{wbII:eq:W:5.17z}
\]
In particular \(\nu^*(\mathcal L_Q|_W)\) is nontrivial.

On every nonzero fibre
\[
 F_t=\mathbb C^2/(Z(s)\mathbb Z^2+\mathbb Z^2),
\]
the first Chern class of \(\mathcal L_Q|_{F_t}\), in the ordered lattice
basis
\((\widehat\gamma,\widehat u,\widehat w,\widehat\delta)\) indexing the
columns of \([Z\mid I_2]\), is exactly
\[
 c_1(\mathcal L_Q|_{F_t})=-2q,
 \qquad
 q=u\wedge w+6\gamma\wedge\delta.
\tag{W.68}\hypertarget{tag-03FW}{}\label{wbII:eq:W:5.17aa}
\]
\end{proposition}



\begin{proof}
\begin{enumerate}[label={$\langle1\rangle$\arabic*.},leftmargin=*]
\item 
\noindent\textbf{Step 1. Local finiteness and Cartier status.}
The fan of \(Y\) consists of the cones over the cells of the locally finite
\(A_2\)-triangulation at height one.  If
\(U_\sigma\simeq\mathbb C^3\) is the affine chart belonging to a triangle
\(\sigma=\{v_0,v_1,v_2\}\), then \(U_\sigma\) meets precisely
\(E_{v_0},E_{v_1},E_{v_2}\).  Indeed, by the orbit--cone
correspondence, \(U_\sigma\cap E_v\ne\varnothing\) exactly when
\(\mathbb R_{\geq0}(v,1)\) is a face of the cone over \(\sigma\).
Thus \(\{E_v\}_{v\in\mathbb Z^2}\) is locally finite and
\(\mathcal D_Q\) is a well-defined locally finite Weil divisor.

Every triangle of the fan is unimodular, as proved in
\CompendiumRef{I}{wb:thm:candidate-cusp-geometry}.  If \(z_0,z_1,z_2\) are the
characters dual to
\((v_0,1),(v_1,1),(v_2,1)\), then
\[
 \mathcal D_Q|_{U_\sigma}
 =\operatorname{div}\!
   \left(z_0^{Q(v_0)}z_1^{Q(v_1)}z_2^{Q(v_2)}\right).
\tag{W.69}\hypertarget{tag-03FX}{}\label{wbII:eq:W:5.17ab}
\]
Negative exponents merely make the displayed monomial meromorphic, as
permitted for a local equation of a Cartier divisor.  Hence
\(\mathcal D_Q\) is Cartier on every chart and therefore on
\(Y_\epsilon\).

\noindent\textbf{Step 2. The signed divisor pullback identity.}
For \(\lambda\in\bar\Lambda=\mathbb Z^2\), set
\[
 b_\lambda=B_0\lambda,
 \qquad
 \ell_\lambda=2Hb_\lambda,
 \qquad
 k_\lambda=Q(b_\lambda),
 \qquad
 m_\lambda=(\ell_\lambda,k_\lambda)\in M'.
\tag{W.70}\hypertarget{tag-03FY}{}\label{wbII:eq:W:5.17ac}
\]
The fan automorphism sends \(E_v\) to \(E_{v+b_\lambda}\).
Consequently \(\Phi_\lambda^*E_w=E_{w-b_\lambda}\), and
\begin{align*}
 \Phi_\lambda^*\mathcal D_Q-\mathcal D_Q
 &=\sum_{v\in\mathbb Z^2}
   \bigl(Q(v+b_\lambda)-Q(v)\bigr)E_v\\
 &=\sum_{v\in\mathbb Z^2}
   \bigl(2v^{\mathsf t}Hb_\lambda+
         b_\lambda^{\mathsf t}Hb_\lambda\bigr)E_v
 =\operatorname{div}(\chi^{m_\lambda}).
\end{align*}
The \(t\)-dependent horizontal torus translation
\(p\mapsto(c_\lambda(t(p)),1)\cdot p\) preserves every
torus-invariant prime divisor \(E_v\).  Hence
\[
 \boxed{\ 
 \Psi_\lambda^*\mathcal D_Q
 =\mathcal D_Q+\operatorname{div}(\chi^{m_\lambda})\ }.
\tag{W.71}\hypertarget{tag-03FZ}{}\label{wbII:eq:W:5.17ad}
\]

\noindent\textbf{Step 3. The complete \(\Psi\)-cocycle.}
The matrices in \eqref{wbII:eq:W:5.17t} and \eqref{wbII:eq:W:5.17x} give
\[
 PC(t)=
 \begin{pmatrix}
 6(b(t)-h(t))&6\mu(t)\\
 6\mu(t)&h(t)
 \end{pmatrix},
\tag{W.72}\hypertarget{tag-03G0}{}\label{wbII:eq:W:5.17ae}
\]
which is symmetric.  Thus
\[
 A(t):=2PC(t)
\tag{W.73}\hypertarget{tag-03G1}{}\label{wbII:eq:W:5.17af}
\]
is a symmetric holomorphic matrix on the full cusp disc.  Define the
holomorphic units \(h_\lambda\) and meromorphic functions \(f_\lambda\) by
\[
 h_\lambda(t)=
 \exp\!\left(\pi i\lambda^{\mathsf t}A(t)\lambda\right),
 \qquad
 f_\lambda=h_\lambda(t)\chi^{m_\lambda}.
\tag{W.74}\hypertarget{tag-03G2}{}\label{wbII:eq:W:5.17ag}
\]
Equation \eqref{wbII:eq:W:5.17ad} gives
\[
 \operatorname{div}(f_\lambda)
 =\Psi_\lambda^*\mathcal D_Q-\mathcal D_Q.
\tag{W.75}\hypertarget{tag-03G3}{}\label{wbII:eq:W:5.17ah}
\]

The cocharacter map underlying \(\Phi_\mu\) is
\[
 \varphi_\mu(y,y_3)=(y+y_3b_\mu,y_3),
\]
so its pullback on characters is
\[
 \varphi_\mu^*(\ell,k)
 =(\ell,k+\ell^{\mathsf t}b_\mu).
\tag{W.76}\hypertarget{tag-03G4}{}\label{wbII:eq:W:5.17ai}
\]
The quadratic identity for \(Q\) gives
\[
 m_{\lambda+\mu}=m_\mu+\varphi_\mu^*m_\lambda.
\tag{W.77}\hypertarget{tag-03G5}{}\label{wbII:eq:W:5.17aj}
\]
Indeed, the first coordinates add, while the last coordinate on the
right is
\[
 Q(b_\mu)+Q(b_\lambda)+2b_\lambda^{\mathsf t}Hb_\mu
 =Q(b_\lambda+b_\mu).
\]
Using \eqref{wbII:eq:W:5.17u}, \eqref{wbII:eq:W:5.17v}, and \eqref{wbII:eq:W:5.17ai}, one obtains
\begin{align*}
 \Psi_\mu^*\chi^{m_\lambda}
 &=\chi^{m_\lambda}(c_\mu(t),1)
   \chi^{\varphi_\mu^*m_\lambda}\\
 &=\exp\!\left(2\pi i\ell_\lambda^{\mathsf t}C(t)\mu\right)
   \chi^{\varphi_\mu^*m_\lambda}\\
 &=\exp\!\left(2\pi i\lambda^{\mathsf t}A(t)\mu\right)
   \chi^{\varphi_\mu^*m_\lambda}.
\tag{W.78}\hypertarget{tag-03G6}{}\label{wbII:eq:W:5.17ak}
\end{align*}
The last equality is exact because
\[
 \ell_\lambda^{\mathsf t}C\mu
 =(2HB_0\lambda)^{\mathsf t}C\mu
 =2\lambda^{\mathsf t}B_0^{\mathsf t}HC\mu
 =2\lambda^{\mathsf t}PC\mu
 =\lambda^{\mathsf t}A\mu.
\]
Since \(A=A^{\mathsf t}\), \eqref{wbII:eq:W:5.17ag}, \eqref{wbII:eq:W:5.17aj}, and \eqref{wbII:eq:W:5.17ak} imply
\begin{align*}
 f_\mu\,\Psi_\mu^*f_\lambda
 &=h_\mu h_\lambda
   \exp(2\pi i\lambda^{\mathsf t}A\mu)
   \chi^{m_\mu+\varphi_\mu^*m_\lambda}\\
 &=\exp\!\left(
   \pi i(\lambda+\mu)^{\mathsf t}A(\lambda+\mu)\right)
   \chi^{m_{\lambda+\mu}}
 =f_{\lambda+\mu}.
\tag{W.79}\hypertarget{tag-03G7}{}\label{wbII:eq:W:5.17al}
\end{align*}

Use the divisor-sheaf convention
\[
 \mathcal O(D)(U)=
 \{g\in\mathcal M_Y(U):
   \operatorname{div}(g)+D|_U\geq0\}.
\]
If \(D'=D+\operatorname{div}(f)\), multiplication by \(f\) maps
\(\mathcal O(D')\) isomorphically to \(\mathcal O(D)\).  Thus
\eqref{wbII:eq:W:5.17ad} and \eqref{wbII:eq:W:5.17ag} define
\[
 \theta_\lambda:
 \Psi_\lambda^*\mathcal O(\mathcal D_Q)
 \xrightarrow{\ \times f_\lambda\ }
 \mathcal O(\mathcal D_Q).
\tag{W.80}\hypertarget{tag-03G8}{}\label{wbII:eq:W:5.17am}
\]
Because
\((\Psi_\lambda\Psi_\mu)^*=\Psi_\mu^*\Psi_\lambda^*\),
equation \eqref{wbII:eq:W:5.17al} is exactly
\[
 \theta_{\lambda+\mu}
 =\theta_\mu\circ\Psi_\mu^*\theta_\lambda.
\tag{W.81}\hypertarget{tag-03G9}{}\label{wbII:eq:W:5.17an}
\]
This proves a complete equivariance datum, not merely invariance of the
divisor class.

\noindent\textbf{Step 4. Analytic descent.}
By \CompendiumRef{I}{wb:thm:candidate-cusp-geometry}, the action on \(Y_\epsilon\) is
free and properly discontinuous.  Every point therefore has an open
neighbourhood \(U\) whose distinct translates are disjoint and on which
\(\pi|_U\) is a biholomorphism onto \(\pi(U)\).  Transport
\(\mathcal O(\mathcal D_Q)|_U\) to \(\pi(U)\), and on overlaps arising
from translates glue by \eqref{wbII:eq:W:5.17am}.  Equation \eqref{wbII:eq:W:5.17an} is precisely the
triple-overlap condition.  The resulting holomorphic line bundle
\(\mathcal L_Q\) on \(N_0\) satisfies
\[
 \pi^*\mathcal L_Q
 \simeq\mathcal O_{Y_\epsilon}(\mathcal D_Q).
\tag{W.82}\hypertarget{tag-03GA}{}\label{wbII:eq:W:5.17ao}
\]

\noindent\textbf{Step 5. The normalized central restriction.}
Let \(j_0:E_0\hookrightarrow Y_\epsilon\).  Since
\(\nu=\pi\circ j_0\), \eqref{wbII:eq:W:5.17ao} gives
\[
 \nu^*(\mathcal L_Q|_W)
 =j_0^*\pi^*\mathcal L_Q
 \simeq\mathcal O_{E_0}(\mathcal D_Q|_{E_0}).
\tag{W.83}\hypertarget{tag-03GB}{}\label{wbII:eq:W:5.17ap}
\]
The only \(E_v\), \(v\ne0\), meeting \(E_0\) are those indexed by the
six \(a_i\) in the statement, and the intersections are the six toric
\((-1)\)-curves \(C_i\).  Since \(Q(0)=0\), there is no
\(E_0|_{E_0}\) term, so
\[
 \mathcal D_Q|_{E_0}
 =\sum_{i=1}^6Q(a_i)C_i,
 \qquad
 (Q(a_1),\ldots,Q(a_6))=(1,-6,-5,1,-6,-5).
\tag{W.84}\hypertarget{tag-03GC}{}\label{wbII:eq:W:5.17aq}
\]
The hexagonal intersection table is
\[
 C_i^2=-1,
 \qquad
 C_i\cdot C_{i\pm1}=1,
 \qquad
 C_i\cdot C_j=0
 \quad(j\notin\{i,i\pm1\}),
\tag{W.85}\hypertarget{tag-03GD}{}\label{wbII:eq:W:5.17ar}
\]
with indices modulo six.  Therefore
\[
 \deg_{C_i}\mathcal O_{E_0}
 \!\left(\sum_jQ(a_j)C_j\right)
 =-Q(a_i)+Q(a_{i-1})+Q(a_{i+1}),
\]
and direct substitution gives
\[
 \begin{array}{c|rrrrrr}
 i&1&2&3&4&5&6\\ \hline
 -Q(a_i)+Q(a_{i-1})+Q(a_{i+1})
 &-12&2&0&-12&2&0.
 \end{array}
\tag{W.86}\hypertarget{tag-03GE}{}\label{wbII:eq:W:5.17as}
\]
The degree on \(C_1=C_{e_1}\) is \(-12\), proving nontriviality.

\noindent\textbf{Step 6. The generic-fibre factor and its exact Chern class.}
Fix \(t\ne0\), choose \(s\) with \(t=e^{2\pi i s}\), and write
\(x=e^{2\pi i\zeta}\), \(\zeta\in\mathbb C^2\).
Equation \eqref{wbII:eq:W:5.17u} identifies the deck transformation indexed by
\(\lambda\in\mathbb Z^2\) with
\[
 \zeta\longmapsto\zeta+Z(s)\lambda
\]
modulo the horizontal translations
\(\zeta\mapsto\zeta+n\), \(n\in\mathbb Z^2\).  The matrix identities
\[
 HB_0=P^{\mathsf t},
 \qquad
 B_0^{\mathsf t}HB_0=PB_0
\tag{W.87}\hypertarget{tag-03GF}{}\label{wbII:eq:W:5.17at}
\]
give
\[
 \chi^{m_\lambda}(x,t)
 =x^{2P^{\mathsf t}\lambda}
  t^{\lambda^{\mathsf t}PB_0\lambda},
\qquad
 h_\lambda(t)
 =\exp(2\pi i\lambda^{\mathsf t}PC(t)\lambda).
\]
Since \(Z=sB_0+C\), it follows that
\[
 f_\lambda(\zeta)
 =\exp\!\left(
 2\pi i\lambda^{\mathsf t}PZ\lambda+
 4\pi i\lambda^{\mathsf t}P\zeta\right).
\tag{W.88}\hypertarget{tag-03GG}{}\label{wbII:eq:W:5.17au}
\]
The matrix \(PZ=sPB_0+PC\) is symmetric:
\[
 PB_0=\operatorname{diag}(-6,1),
\]
and \(PC\) is symmetric by \eqref{wbII:eq:W:5.17ae}.

Under the descent convention \eqref{wbII:eq:W:5.17am}, a pulled-back local section
written as \(\varphi(\zeta)\) in the canonical frame satisfies
\[
 \varphi(\zeta+Z\lambda)
 =f_\lambda(\zeta)^{-1}\varphi(\zeta).
\]
Thus the factor of automorphy for the lattice translation
\(\zeta\mapsto\zeta+Z\lambda+n\) is
\[
 a_{(\lambda,n)}(\zeta)
 =\exp(2\pi i g_{(\lambda,n)}(\zeta)),
\qquad
 g_{(\lambda,n)}(\zeta)
 =-\lambda^{\mathsf t}PZ\lambda
  -2\lambda^{\mathsf t}P\zeta.
\tag{W.89}\hypertarget{tag-03GH}{}\label{wbII:eq:W:5.17av}
\]
This is independent of \(n\) as a factor of automorphy: replacing
\(\zeta\) by \(\zeta+n\) changes \(g\) by the integer
\(-2\lambda^{\mathsf t}Pn\).

For \(l=(\lambda,n)\), \(l'=(\mu,m)\), the exponential connecting
cocycle is
\begin{align*}
 c(l,l')
 &=g_l(\zeta+Z\mu+m)+g_{l'}(\zeta)-g_{l+l'}(\zeta)\\
 &=-2\lambda^{\mathsf t}Pm\in\mathbb Z,
\tag{W.90}\hypertarget{tag-03GI}{}\label{wbII:eq:W:5.17aw}
\end{align*}
where the symmetry of \(PZ\) cancels the two mixed \(Z\)-terms.
The corresponding alternating integral form is
\[
 E(l,l')=c(l,l')-c(l',l)
 =-2\lambda^{\mathsf t}Pm+2\mu^{\mathsf t}Pn.
\tag{W.91}\hypertarget{tag-03GJ}{}\label{wbII:eq:W:5.17ax}
\]
This is \(c_1(\mathcal L_Q|_{F_t})\).  In the ordered column basis of
\([Z\mid I_2]\), the cross-block of the class
\(q=u\wedge w+6\gamma\wedge\delta\) from
\((\widehat\gamma,\widehat u)\) to
\((\widehat w,\widehat\delta)\) is
\[
 \begin{pmatrix}
 q(\widehat\gamma,\widehat w)&
 q(\widehat\gamma,\widehat\delta)\\
 q(\widehat u,\widehat w)&
 q(\widehat u,\widehat\delta)
 \end{pmatrix}
 =\begin{pmatrix}0&6\\1&0\end{pmatrix}=P.
\tag{W.92}\hypertarget{tag-03GK}{}\label{wbII:eq:W:5.17ay}
\]
Hence \eqref{wbII:eq:W:5.17ax} is exactly \(-2q\), as asserted.

\end{enumerate}
\end{proof}